\documentclass[11pt,a4paper]{article}
\usepackage[T1]{fontenc}
\usepackage{lmodern}
\usepackage[margin=2.3cm]{geometry}
\usepackage{amsmath,amssymb}
\usepackage{booktabs}
\usepackage{microtype}
\usepackage{caption}
\usepackage{pgfplots}
\pgfplotsset{compat=1.17}
\usepgfplotslibrary{groupplots}
\usepackage[colorlinks=true,linkcolor=blue!50!black,urlcolor=blue!50!black]{hyperref}
\setlength{\emergencystretch}{3em}

\newcommand{\D}{\mathcal{D}}
\newcommand{\Z}{\mathcal{Z}}
\newcommand{\C}{\mathcal{C}}
\newcommand{\T}{\mathcal{T}}
\newcommand{\pp}{{\uparrow\uparrow}}
\newcommand{\pa}{{\uparrow\downarrow}}
\newcommand{\cin}{c_e^{\rm in}}
\newcommand{\DIR}{\mathrm{DIR}}
\newcommand{\RLC}{\mathrm{RLC}}
\newcommand{\BL}{\mathrm{BL}}

\definecolor{parcol}{HTML}{A32D2D}
\definecolor{antcol}{HTML}{2A78D6}

\title{CQD analysis of the SG0--SG1 spin-resolved beam separations\\[4pt]
\large Data set of 2026-07-06: derivation and fits}
\author{Companion to \emph{Analytical Model and Provisional Calibration of the}\\
\emph{Null Stern--Gerlach CQD Experiment}, Revision F (October 13, 2026)}
\date{October 16, 2026}

\begin{document}
\maketitle

\begin{abstract}
\noindent
The double-stage co-quantum-dynamics (CQD) model of the source note is fitted to the spin-resolved
beam-separation data of 2026-07-06 (\texttt{spin\_resolved\_sep\_20260706.jld2}): 15 SG0 currents
from 0 to 3.64~A for parallel and antiparallel SG0--SG1 fields, each with its own uncertainty. The
chain of equations from the oven ensemble to the detector separation is derived in full, using the
notation of the source note; equation numbers in square brackets refer to it. The SG0 field is
divided by $f=7$ (Revision~F scaling); $f$ is a fixed calibration input, not a fitting parameter.
In the recommended reduction the parallel set alone is fitted with pure CQD, which gives
$k_i=(2.67^{+0.60}_{-0.48})\times10^{-6}$ ($\chi^2=2.8$ for 15 points). With this $k_i$ fixed, pure
CQD predicts the wrong sign for the antiparallel set ($\chi^2=109$); a retained-label fraction
proportional to the SG0 dose, $w_{\rm ret}\simeq0.036~\text{A}^{-1}\times I_0$, reproduces it
($\chi^2=0.6$). With the measured SG0 field ($f=1$) the same procedure gives
$k_i=(1.52\pm0.09)\times10^{-7}$. The alternative fitting strategies are compared, and the
dependence on $f$ is mapped. Fitted curves are supplied as CSV files.
\end{abstract}

\tableofcontents

%==========================================================================
\section{The data}
\label{sec:data}

The beam passes first through SG0, a uniform-field magnet with no gradient, and then through SG1, a
gradient Stern--Gerlach magnet at the fixed current $I_1=40$~mA. The camera records the distance
between the two beam peaks. The data file contains, for SG0 and SG1 fields parallel ($\pp$, group
\texttt{idx\_7}) and antiparallel ($\pa$, group \texttt{idx\_8}), the change of this separation
relative to its zero-current value, $\Delta D(I_0)$ in pixels, together with a standard uncertainty for
every point (Table~\ref{tab:points}, Figure~\ref{fig:data}).

Compared with the digitized eight-point sets used previously, the new data
\begin{itemize}
\item contain 15 currents per configuration, including five low currents $0\le I_0\le0.125$~A and two
near-duplicate pairs (1.403/1.404~A and 3.598/3.642~A);
\item carry individual uncertainties $\sigma_j=0.47$--$0.53$~px, used as weights in every fit;
\item have small zero-current values, $-0.047$~px (parallel) and $+0.049$~px (antiparallel). Their
difference, 0.096~px, equals the stored quantity \texttt{y0}. The model has $\Delta D(0)=0$ by
construction; the offsets are a tenth of $\sigma$ and are kept as data.
\end{itemize}
The model normalizes all separations to the absolute zero-current separation, which is not stored in
the file. The value $D_0=13$~pixels of the source note is used throughout.

\begin{figure}[h]
\centering
\begin{tikzpicture}
\begin{axis}[width=0.86\textwidth,height=0.48\textwidth,xlabel={SG0 current $I_0$ (A)},
  ylabel={$\Delta D$ (px)},xmin=-0.1,xmax=3.8,ymin=-2.8,ymax=5.6,grid=major,grid style={gray!20},
  legend pos=north west,legend style={font=\footnotesize,draw=none,fill=white,fill opacity=0.85,text opacity=1},legend cell align=left]
\addplot[only marks,mark=diamond*,mark size=3pt,parcol,error bars/.cd,y dir=both,y explicit]
  table[x=I0_A,y=dD_parallel_px,y error=sigma_parallel_px,col sep=comma]{spin_resolved_sep_20260706.csv};
\addlegendentry{parallel $\pp$}
\addplot[only marks,mark=square*,mark size=2pt,antcol,error bars/.cd,y dir=both,y explicit]
  table[x=I0_A,y=dD_antiparallel_px,y error=sigma_antiparallel_px,col sep=comma]{spin_resolved_sep_20260706.csv};
\addlegendentry{antiparallel $\pa$}
\addplot[gray,dashed,forget plot] coordinates{(-0.1,0)(3.8,0)};
\end{axis}
\end{tikzpicture}
\caption{The 2026-07-06 data: change of the beam separation relative to zero SG0 current.}
\label{fig:data}
\end{figure}

\paragraph{Physical content.} A uniform field exerts no force, and standard quantum mechanics with the
full hyperfine Hamiltonian predicts no change of the separation (section~23 of the source note). In
CQD, SG0 \emph{rotates the electron magnetic moment} towards one pole without deflecting the atom.
This changes how strongly each atom is pushed inside SG1, and therefore changes the separation. The
size and sign of that change test the CQD collapse law and the electron--nuclear branching rule.

%==========================================================================

\section{Model derivation}

\subsection{Initial ensemble and heart-shaped branching}

At the oven exit the electron and nuclear moments are independent and isotropic. In the cosine
variables $c_e=\cos\theta_e$ and $c_n=\cos\theta_n$ the measure is uniform on a square [89]:
\begin{equation}
dP_0=\tfrac14\,dc_e\,dc_n,\qquad -1\le c_e,c_n\le 1.
\end{equation}
CQD assigns the electron branch by comparing the two polar angles [90]:
\begin{equation}
q=\operatorname{sgn}(\theta_n-\theta_e)=\operatorname{sgn}(c_e-c_n).
\end{equation}
The $q=+1$ branch is the triangle $c_n<c_e$. Integrating $c_n$ over $-1<c_n<c_e$ gives the length
$1+c_e$, so that
\begin{equation}
p(c_e\mid q{=}+1)=\frac{1+c_e}{2},\qquad p(c_n\mid q{=}+1)=\frac{1-c_n}{2},\qquad
\langle c_e\rangle_+=\tfrac13 .
\end{equation}
The nuclear conditional density $(1-\cos\theta_n)/4\pi$ is the cardioid (``heart'') of the source
note. In the rest of the derivation the heart shape appears only through the weight $(1+c_e)$, which
is fixed by geometry and is not a fitting weight. It also explains why the weak-field response of a
branch is $1/3$ rather than $0$ or $1$.

\subsection{CQD collapse as a translation}

The CQD collapse law [118] is
\begin{equation}
\tan\frac{\theta_e(t)}2=\tan\frac{\theta_e(0)}2\,e^{-q\D_c(t)},\qquad \D_c=k_i\,|\Delta\phi_e|,
\end{equation}
where $\D_c$, the \emph{collapse dose}, is $k_i$ times the accumulated Larmor phase. With the electron
log-amplitude ratio
\begin{equation}
\C_e=\ln\frac{|A_\uparrow|}{|A_\downarrow|}=-\ln\tan\frac{\theta_e}{2}=\operatorname{artanh}c_e,
\qquad c_e=\tanh\C_e,
\end{equation}
the law becomes a uniform translation [129]:
\begin{equation}
\boxed{\;\C_e^{\rm out}=\C_e^{\rm in}+q\,\D_c\;}
\end{equation}
Doses of successive stages with the same branch therefore add. For a stage with longitudinal field
profile $B(y)$ traversed at speed $v_a$ the dose depends only on the field integral [131]:
\begin{equation}
\D_c=\frac{k_i|\gamma_e|}{v_a}\int|B(y)|\,dy=\frac{k_i|\gamma_e|\,\BL}{v_a}.
\end{equation}
Converting back to a cosine with the addition theorem for $\tanh$ gives the transfer map [145]:
\begin{equation}
\T_q(c;\D)=\frac{c+q\tanh\D}{1+q\,c\tanh\D}\;=\;c+q(1-c^2)\D+O(\D^2).
\end{equation}
The factor $1-c^2$ is a logistic flow: the moment turns fastest at the equator and stops at the poles.

\subsection{Linking the SG0 and SG1 doses}

The constant $k_i|\gamma_e|/v_a$ cancels in the ratio of the two doses [138--144]:
\begin{equation}
\frac{\D_{c,0}(I_0)}{\D_{c,1}}=\frac{\BL_0(I_0)}{\BL_1}
\simeq\frac{\beta_0I_0\,\ell_0^{\rm eff}}{B_1\,\ell_1^{\rm eff}}
=\frac{I_0}{I_\star},\qquad
I_\star=\frac{B_1\ell_1^{\rm eff}}{\beta_0\ell_0^{\rm eff}} .
\label{eq:link}
\end{equation}
With $\beta_0=0.031895$~T\,A$^{-1}$, $B_1=0.018931$~T, $\ell_0^{\rm eff}=0.040$~m and
$\ell_1^{\rm eff}=0.070$~m one obtains $I_\star=1.03869$~A. Thus $\D_{c,0}=\alpha_0I_0$ with
$\alpha_0=\D_{c,1}/I_\star$. \emph{Once $\D_{c,1}$ is chosen, the SG0 dose at every current is fixed};
SG0 introduces no free parameter.

\subsection{Interstage transfer}

After SG0 each atom has $c_{e,0}^{\rm out}=\T_{q_0}(\cin;\D_{c,0})$.

\paragraph{Parallel fields.} The field direction does not change, the electron keeps its branch,
$q_1=q_0$, and enters SG1 pre-aligned towards the pole that SG1 also drives it to [155--156]:
$c_{e,1}^{\rm in}=\T_{q_0}(\cin;\D_{c,0})$.

\paragraph{Antiparallel fields.} The field reverses adiabatically; the electron follows it while the
slow nucleus does not [157--165]:
\begin{equation}
\theta_e\to\pi-\theta_e\ (c_e\to-c_e),\qquad \theta_n\to\theta_n,\qquad
q_1^\pa=-\operatorname{sgn}\!\big(c_{e,0}^{\rm out}+c_n\big).
\end{equation}
The SG1 branch boundary is $\theta_n+\theta_{e,0}^{\rm out}=\pi$. Each SG0 triangle is cut into two
classes: \emph{reversed-label} atoms ($q_1=-q_0$), which continue towards the new pole and are well
prepared, and \emph{retained-label} atoms ($q_1=q_0$), which SG1 must first drag back through the
equator (counter-collapse) and which receive a weak, initially negative, push.

\subsection{Force in SG1 and displacement at the detector}

Inside SG1, $\C_e(t)=\C_{e,1}^{\rm in}+q_1r_1t$ with $r_1=k_i|\gamma_e|B_1$, and the force is
[187--190]
\begin{equation}
F_z(t)=\mu_eG_1\tanh\!\big(\C_{e,1}^{\rm in}+q_1r_1t\big).
\end{equation}
A momentum increment at time $t$ acts for the rest of the magnet and the whole drift, so the detector
displacement is [194]
\begin{equation}
z_d=\frac{\mu_eG_1}{M_a}\int_0^{T_1}(T_1+T_d-t)\,\tanh\!\big(\C_{e,1}^{\rm in}+q_1r_1t\big)\,dt,
\qquad T_1=\frac{\ell_1^{\rm eff}}{v_a},\quad T_d=\frac{L_d}{v_a}.
\end{equation}
Dividing by the fully aligned value and using $\lambda=r_1t$, $\D_{c,1}=r_1T_1$ and
$\rho=r_1T_d=\DIR\cdot\D_{c,1}$ with $\DIR=T_d/T_1$ gives the normalized per-atom response
[204, 211]:
\begin{equation}
\boxed{\;
\tilde z_{\rm ex}(c;\D_{c,1},\rho)=
\frac{\displaystyle\int_0^{\D_{c,1}}(\D_{c,1}+\rho-\lambda)\tanh(\operatorname{artanh}c+\lambda)\,d\lambda}
{\rho\,\D_{c,1}+\D_{c,1}^2/2}\;}
\label{eq:zex}
\end{equation}
It lies between $-1$ and $+1$. In the long-drift limit it reduces to
$\tilde z=\D_{c,1}^{-1}\ln(\cosh\D_{c,1}+c\sinh\D_{c,1})$, and for a small dose
$\tilde z\simeq c+\tfrac12\D_{c,1}(1-c^2)$: each atom is pushed in proportion to its average
$c_e$ during the passage. The closed form in terms of the dilogarithm [203--207] is the analytic
evaluation of the same integral.

\subsection{Ensemble responses}

By inversion symmetry the normalized separation equals the $q_1=+1$ branch centroid [213--216].
Writing $\tilde z(\cdot)\equiv\tilde z_{\rm ex}(\cdot;\D_{c,1},\DIR\,\D_{c,1})$:
\begin{align}
\Z^\pp&=\tfrac12\int_{-1}^{1}(1+c)\,\tilde z\big(\T_+(c;\D_{c,0})\big)\,dc
&&\text{[217]}\label{eq:Zpar}\\
\Z^\pa_{\rm geom}&=\tfrac12\int_{-1}^{1}\Big[L_+(c)\,\tilde z\big(-\T_+(c)\big)
+L_-(c)\,\tilde z\big(-\T_-(c)\big)\Big]dc
&&\text{[218--221]}\\
&\qquad L_+=1+\min[c,-\T_+(c)],\quad L_-=\max[0,-\T_-(c)-c]\nonumber\\
\Z^\pa_{\rm ret}&=\tfrac12\int_{-1}^{1}(1+c)\,\tilde z\big(-\T_+(c)\big)\,dc
&&\text{[251]}
\end{align}
$L_\pm$ are the heart-shaped weights cut by the new SG1 branch boundary. $\Z_{\rm ret}$ assumes that
every atom keeps its label. The antiparallel response actually fitted is the mixture [252--253]
\begin{equation}
\Z^\pa=(1-w_{\rm ret})\,\Z^\pa_{\rm geom}+w_{\rm ret}\,\Z^\pa_{\rm ret},\qquad
w_{\rm ret}=1-e^{-\RLC\,\D_{c,0}} .
\end{equation}
The retained-label term is phenomenological. It stands in for unresolved fringe-field dynamics and is
not part of pure CQD.

\subsection{Absolute scale from the 13-pixel baseline}

At $I_0=0$ both configurations reduce to $\Z_0(\D_{c,1})=\Z^\pp(0,\D_{c,1})$, with
$\Z_0\simeq\tfrac13+\tfrac13\D_{c,1}$ for a weak dose. Because $D=D_{\rm sat}\Z$ and $D(0)=13$~px,
the mechanical scale cancels and the model curve is [240--241]
\begin{equation}
\boxed{\;
\Delta D_s(I_0)=13\left[\frac{\Z_s\big(\D_{c,1}I_0/I_\star,\;\D_{c,1}\big)}{\Z_0(\D_{c,1})}-1\right],
\qquad s=\pp,\pa .\;}
\label{eq:model}
\end{equation}
The remaining geometric quantity, $\DIR$, is eliminated with a second use of the baseline [249]:
\begin{equation}
13=K_{\rm pix}\big(\DIR+\tfrac12\big)\,\Z_{0,\rm ex}(\D_{c,1},\DIR),\qquad
K_{\rm pix}=\frac{2\mu_eG_1T_1^2}{M_ad_{\rm pix}}=6.47~\text{px}.
\label{eq:dir}
\end{equation}
For every trial $\D_{c,1}$ the scalar equation~\eqref{eq:dir} is solved for $\DIR$. The drift length
$L_d=\DIR\cdot\ell_1^{\rm eff}$ is therefore an \emph{inferred} quantity.

%==========================================================================
\subsection{Leading-order picture of the parallel curve}
\label{sec:leading}

Expanding to first order in both doses, $\T_+(c)\simeq c+\D_{c,0}(1-c^2)$ and
$\tilde z\simeq c+\tfrac12\D_{c,1}(1-c^2)$. The heart-weighted integrals are
$\tfrac12\int(1+c)c\,dc=\tfrac13$ and $\tfrac12\int(1+c)(1-c^2)\,dc=\tfrac23$, hence
\begin{equation}
\Z^\pp\simeq\tfrac13+\tfrac23\D_{c,0}+\tfrac13\D_{c,1}
\quad\Longrightarrow\quad
\frac{\Delta D^\pp}{13}\simeq2\,\D_{c,0}=2\alpha_0I_0 .
\end{equation}
With the $f=1$ value $\alpha_0=0.057$~A$^{-1}$ this predicts about $1.5$~px\,A$^{-1}$, against about $1.3$~px\,A$^{-1}$ in the data; the difference
comes from the $\D_{c,1}$ term in $\Z_0$ and from the onset of saturation. \emph{The slope of the
parallel curve measures, almost directly, how fast SG0 rotates electron moments per ampere.} This is
why $\alpha_0$, and hence $k_i$, is robust against the SG1 calibration, while $\D_{c,1}$ is not.

%==========================================================================
\section{Calibration inputs and fitting method}
\label{sec:method}

\paragraph{SG0 field scaling.} Revision~F divides the SG0 field, and hence $\BL_0$, by $f=7$. Only the
dose link~\eqref{eq:link} changes:
\begin{equation}
\D_{c,0}(I_0)=\D_{c,1}\,\frac{I_0}{f\,I_\star},\qquad fI_\star=7\times1.03869~\text{A}=7.2708~\text{A}.
\end{equation}
Because $\BL_0$ enters only through $B_0\ell_0^{\rm eff}$, this is equivalent to an effective SG0 length
$\ell_0^{\rm eff}=0.040/7=0.571$~cm at the measured field. $f$ is a \emph{fixed calibration input}; the
measured case $f=1$ is reported for comparison and Section~\ref{sec:fscan} maps the dependence on $f$.

\paragraph{Fixed inputs.} $\beta_0=0.031895$~T\,A$^{-1}$, $B_1=0.018931$~T, $\ell_0^{\rm eff}=0.040$~m,
$\ell_1^{\rm eff}=0.070$~m, $G_1=10.779$~T\,m$^{-1}$, $v_a=600.5$~m\,s$^{-1}$,
$|\gamma_e|=1.76086\times10^{11}$~rad\,s$^{-1}$\,T$^{-1}$, $K_{\rm pix}=6.47$~px, $D_0=13$~px. Unless
stated otherwise, $\DIR$ is solved from the baseline equation~\eqref{eq:dir}.

\paragraph{Objective.} Each data set is fitted by minimizing
\begin{equation}
\chi^2=\sum_{j=1}^{15}\left(\frac{\Delta D_j^{\rm model}-\Delta D_j^{\rm data}}{\sigma_j}\right)^2
\end{equation}
with the per-point uncertainties of the data file. Uncertainties on fitted parameters are profile
$\Delta\chi^2=1$ intervals (one-parameter fits) or curvature errors $\sqrt{2/\chi''}$. The model is
evaluated with Gauss--Legendre quadrature of \eqref{eq:zex}--\eqref{eq:Zpar}, and $\DIR$ by bisection
of~\eqref{eq:dir}.

\paragraph{Free parameters.} The parallel model has one free parameter, $\D_{c,1}$, equivalent to
$k_i=\D_{c,1}v_a/(|\gamma_e|\BL_1)$. Pure CQD adds none for the antiparallel set, so its antiparallel
curve is a prediction once $k_i$ is fixed. The phenomenological retained-label term adds one
parameter, $\kappa$ or $\RLC$.

%==========================================================================
\section{Recommended fit at \texorpdfstring{$f=7$}{f=7}: parallel set first, then \texorpdfstring{$k_i$}{ki} borrowed}
\label{sec:recommended}

\paragraph{Step 1: parallel set, pure CQD.} Minimizing $\chi^2_\pp(\D_{c,1})$ gives
\begin{equation}
\D_{c,1}=1.037,\qquad k_i=\frac{\D_{c,1}v_a}{|\gamma_e|\BL_1}=(2.67^{+0.60}_{-0.48})\times10^{-6},\qquad
\chi^2_\pp=2.83~(15~\text{points}),
\end{equation}
with $\alpha_0=\D_{c,1}/(fI_\star)=0.143$~A$^{-1}$, $\D_{c,0}(3.64~\text{A})=0.52$, $\Z_{0,\rm ex}=0.603$
and $\DIR=2.83$ ($L_d=0.198$~m, inferred). Fixing instead the reported $L_d=0.403$~m ($\DIR=5.757$)
gives $k_i=(2.59^{+0.51}_{-0.43})\times10^{-6}$ with $\chi^2_\pp=4.10$.

\paragraph{Step 2: antiparallel set with $k_i$ fixed.} Pure CQD now has no free parameter. Its strict
mirror-and-rebranch prediction rises to $+0.82$~px at 3.64~A, whereas the data fall to $-2.01$~px:
$\chi^2_\pa=109$. Fitting only the retained-label weight,
\begin{equation}
\Z^\pa=(1-w_{\rm ret})\Z^\pa_{\rm geom}+w_{\rm ret}\Z^\pa_{\rm ret},\qquad
w_{\rm ret}=\kappa\,\D_{c,0}\quad\text{or}\quad w_{\rm ret}=1-e^{-\RLC\,\D_{c,0}},
\end{equation}
gives $\kappa=0.249\pm0.024$ ($\chi^2_\pa=0.62$) or $\RLC=0.264\pm0.027$ ($\chi^2_\pa=0.80$). Both
forms vanish at $\D_{c,0}=0$ and so preserve the common baseline.

\begin{center}
\fbox{\parbox{0.92\textwidth}{\centering\textbf{Key result ($f=7$ fixed, $\DIR$ from the 13-px baseline, 2026-07-06 data)}\\[4pt]
Parallel, pure CQD: $\;k_i=(2.67^{+0.60}_{-0.48})\times10^{-6}$ ($\D_{c,1}=1.037$), $\;\chi^2=2.83$ for 15 points.\\[2pt]
Antiparallel, $k_i$ fixed: pure CQD $\chi^2=109$; with $w_{\rm ret}=\kappa\D_{c,0}$,
$\kappa=0.249\pm0.024$, $\chi^2=0.62$ for 15 points.\\[2pt]
Retained fraction: $w_{\rm ret}=\kappa\alpha_0I_0\approx0.036~\text{A}^{-1}\times I_0$, i.e.\ 13\% at 3.64~A.}}
\end{center}

\begin{figure}[h]
\centering
\begin{tikzpicture}
\begin{axis}[width=0.86\textwidth,height=0.52\textwidth,xlabel={SG0 current $I_0$ (A)},
  ylabel={$\Delta D$ (px)},xmin=-0.1,xmax=3.8,ymin=-3.0,ymax=6.6,grid=major,grid style={gray!20},
  legend pos=north west,legend style={font=\footnotesize,draw=none,fill=white,fill opacity=0.85,text opacity=1},legend cell align=left]
\addplot[parcol,thick] table[x=I0_A,y=dD_parallel_px,col sep=comma]{fit_curves_f7_20260706.csv};
\addlegendentry{parallel: pure CQD fit, $k_i=2.67\times10^{-6}$}
\addplot[antcol,thick] table[x=I0_A,y=dD_antiparallel_px,col sep=comma]{fit_curves_f7_20260706.csv};
\addlegendentry{antiparallel: same $k_i$, $w_{\rm ret}=\kappa\D_{c,0}$}
\addplot[gray,densely dotted,thick] table[x=I0_A,y=dD_antiparallel_pureCQD_px,col sep=comma]{fit_curves_f7_20260706.csv};
\addlegendentry{antiparallel: same $k_i$, pure CQD}
\addplot[only marks,mark=diamond*,mark size=3pt,parcol,error bars/.cd,y dir=both,y explicit]
  table[x=I0_A,y=dD_parallel_px,y error=sigma_parallel_px,col sep=comma]{spin_resolved_sep_20260706.csv};
\addlegendentry{parallel data}
\addplot[only marks,mark=square*,mark size=2pt,antcol,error bars/.cd,y dir=both,y explicit]
  table[x=I0_A,y=dD_antiparallel_px,y error=sigma_antiparallel_px,col sep=comma]{spin_resolved_sep_20260706.csv};
\addlegendentry{antiparallel data}
\addplot[gray,dashed,forget plot] coordinates{(-0.1,0)(3.8,0)};
\end{axis}
\end{tikzpicture}
\caption{Recommended reduction at $f=7$. The parallel curve is a one-parameter pure CQD fit; its $k_i$
is used unchanged for the antiparallel set, where only the retained-label weight is fitted. Curves are
tabulated in \texttt{fit\_curves\_f7\_20260706.csv}.}
\label{fig:recommended}
\end{figure}

\begin{table}[h]
\centering
\caption{Data and fit values at the measured currents ($\Delta D$ in px). ``$f=7$'' is the recommended
fit; ``$f=1$'' is the parallel fit with the measured SG0 field. The last column is the pure CQD
antiparallel prediction at the $f=7$ value of $k_i$.}
\label{tab:points}
\footnotesize\setlength{\tabcolsep}{5pt}
\begin{tabular}{lcccccc}
\toprule
 & \multicolumn{3}{c}{parallel} & \multicolumn{3}{c}{antiparallel}\\
\cmidrule(lr){2-4}\cmidrule(lr){5-7}
$I_0$ (A) & data & fit $f=7$ & fit $f=1$ & data & fit $f=7$ & pure CQD\\
\midrule
0.000 & $-0.047\pm0.50$ & 0.000 & 0.000 & $0.049\pm0.48$ & 0.000 & 0.000 \\
0.005 & $-0.045\pm0.50$ & 0.007 & 0.007 & $0.050\pm0.48$ & $-$0.002 & 0.000 \\
0.012 & $-0.043\pm0.50$ & 0.018 & 0.017 & $0.052\pm0.48$ & $-$0.005 & 0.000 \\
0.043 & $-0.030\pm0.51$ & 0.064 & 0.060 & $0.058\pm0.49$ & $-$0.019 & 0.000 \\
0.125 & $0.028\pm0.51$ & 0.185 & 0.173 & $0.063\pm0.50$ & $-$0.057 & 0.001 \\
0.431 & $0.442\pm0.50$ & 0.628 & 0.594 & $-0.016\pm0.52$ & $-$0.197 & 0.012 \\
0.937 & $1.343\pm0.47$ & 1.325 & 1.282 & $-0.293\pm0.51$ & $-$0.436 & 0.056 \\
1.404 & $2.034\pm0.47$ & 1.930 & 1.909 & $-0.551\pm0.50$ & $-$0.664 & 0.126 \\
1.403 & $2.033\pm0.47$ & 1.929 & 1.908 & $-0.551\pm0.50$ & $-$0.664 & 0.126 \\
1.867 & $2.539\pm0.50$ & 2.493 & 2.522 & $-0.780\pm0.51$ & $-$0.900 & 0.222 \\
2.327 & $3.014\pm0.53$ & 3.018 & 3.124 & $-1.033\pm0.52$ & $-$1.146 & 0.343 \\
2.780 & $3.608\pm0.52$ & 3.500 & 3.707 & $-1.357\pm0.53$ & $-$1.399 & 0.485 \\
3.221 & $4.296\pm0.49$ & 3.939 & 4.267 & $-1.729\pm0.52$ & $-$1.659 & 0.645 \\
3.642 & $4.775\pm0.47$ & 4.329 & 4.793 & $-2.011\pm0.51$ & $-$1.920 & 0.815 \\
3.598 & $4.748\pm0.47$ & 4.290 & 4.738 & $-1.991\pm0.51$ & $-$1.892 & 0.797 \\
\bottomrule
\end{tabular}
\end{table}

\paragraph{Residuals.} With $f=7$ the parallel curve bends over: the model lies 0.45~px below the
3.64~A point (about $1\sigma$) and 0.15--0.19~px above the data at 0.125 and 0.431~A. With $f=1$ the
parallel curve is nearly straight and all residuals are below 0.16~px ($\chi^2=0.48$). The
low-current points lie systematically below both models (the data start more slowly than linearly);
the effect is a third of $\sigma$ and is not significant with the present uncertainties.

\paragraph{Goodness of fit.} All $\chi^2$ values are well below the number of points (reduced
$\chi^2\approx0.04$--0.2). The quoted uncertainties are therefore conservative relative to the scatter,
or the points are correlated (for example through a common baseline subtraction, or if the closely
spaced low-current points are not independent measurements). The $\Delta\chi^2=1$ intervals should be
read with this in mind.

\subsection{Interpretation}
\begin{itemize}
\item \emph{$k_i$ is set by the slope of the parallel curve.} To leading order
$\Delta D^\pp/13\simeq2\D_{c,0}=2\alpha_0I_0$ (Section~\ref{sec:leading}), so the parallel data measure
$\alpha_0$, the SG0 dose per ampere. Converting $\alpha_0$ into $k_i$ uses the calibration, above all
$f$.
\item \emph{The factor 7 does not propagate as a factor 7.} With $f=7$ one needs $\D_{c,1}\approx1$,
where SG1 alone already aligns most of the branch ($\Z_0$ rises from 0.35 to 0.60). The relative effect
of a given SG0 pre-alignment is then smaller, and a larger SG0 dose is needed: $\alpha_0$ rises from
0.057 to 0.143~A$^{-1}$, and $k_i$ by a factor 17.5 rather than 7.
\item \emph{A small retained fraction explains the antiparallel decrease.} Retained-label atoms enter
SG1 counter-aligned and push strongly in the wrong direction, so 13\% at 3.64~A suffices. This is a
one-parameter phenomenological description, not a CQD prediction; whether it is physical must come
from a model of the field-reversal region (transit time $t_f$, adiabaticity, nuclear-angle motion).
\end{itemize}

\subsection{Results with the measured SG0 field (\texorpdfstring{$f=1$}{f=1})}
The same two-step procedure with $f=1$ gives $\D_{c,1}=0.0592$, $k_i=(1.52\pm0.09)\times10^{-7}$
($\chi^2_\pp=0.48$), $\DIR=5.20$ ($L_d=0.364$~m). For the antiparallel set the strict rule gives
$\chi^2=128$, and the retained-label weight gives $\kappa=0.452\pm0.040$ ($\chi^2=0.80$) or
$\RLC=0.470\pm0.044$ ($\chi^2=0.98$), i.e.\ $w_{\rm ret}(3.64~\text{A})=0.094$. Fixing $L_d=0.403$~m
leaves $k_i$ unchanged to four digits. These values agree with the digitized-data analysis
($1.53\times10^{-7}$, $\RLC=0.466$).

%==========================================================================
\section{Comparison of fitting strategies}
\label{sec:strategies}

\begin{table}[h]
\centering
\caption{Fitting strategies applied to the 2026-07-06 data. $\chi^2$ is for 15 points per set.}
\label{tab:strategies}
\footnotesize\setlength{\tabcolsep}{4pt}
\begin{tabular}{p{3.5cm}p{2.0cm}p{3.6cm}p{5.2cm}}
\toprule
Strategy & Free & $k_i$ ($f=1$ / $f=7$) & Verdict\\
\midrule
1. Both sets together (Figure-12 method) & $\D_{c,1}$, $\RLC$ & $1.52\times10^{-7}$ / $2.68\times10^{-6}$ ($\chi^2=1.46$ / 3.63, 30 pts) &
Works, but the non-CQD term $\RLC$ takes part in the fit that defines $k_i$. Gives the same $k_i$ as
strategy~2 within 0.5\%.\\[3pt]
2. Parallel only, then $k_i$ borrowed & $\D_{c,1}$; then $\kappa$ (or $\RLC$) & $1.52\times10^{-7}$ / $2.67\times10^{-6}$ ($\chi^2_\pp=0.48$ / 2.83) &
\textbf{Recommended.} $k_i$ comes from pure CQD on the data pure CQD can fit; the antiparallel set
then tests the reversal physics.\\[3pt]
3a. Antiparallel only, pure CQD & $\D_{c,1}$ & none &
Fails: the best it can do is $k_i\to0$, i.e.\ no change ($\chi^2\approx58$--60).\\[3pt]
3b. Antiparallel only, with retained label & $\D_{c,1}$, $\RLC$ & undetermined &
$\RLC$ compensates any $k_i$; $\chi^2$ stays near 1 (Figure~\ref{fig:strategies}, left).\\[3pt]
4. $f$ free as well & $\D_{c,1}$, ($\RLC$), $f$ & undetermined in $1.5\times10^{-7}$--$3\times10^{-6}$ &
Only $\alpha_0$ is measured; $\chi^2$ is flat for $f\lesssim6$ (Section~\ref{sec:fscan}).\\
\bottomrule
\end{tabular}
\end{table}

\begin{figure}[h]
\centering
\begin{tikzpicture}
\begin{axis}[width=0.62\textwidth,height=0.4\textwidth,xlabel={$k_i$},ylabel={$\chi^2_\pa$ (15 points)},
  xmode=log,ymode=log,xmin=5e-9,xmax=1e-5,ymin=0.1,ymax=4000,grid=major,grid style={gray!20},
  legend pos=north west,legend style={font=\scriptsize,draw=none,fill=white,fill opacity=0.85,text opacity=1},legend cell align=left]
\addplot[black,thick,densely dotted,mark=*,mark size=1.2pt] table[x=ki,y=chi2_f1_pure,col sep=comma]{plot_strategy3.csv};
\addlegendentry{$f=1$, pure CQD}
\addplot[black,thick,mark=*,mark size=1.2pt] table[x=ki,y=chi2_f1_rlc,col sep=comma]{plot_strategy3.csv};
\addlegendentry{$f=1$, with $\RLC$}
\addplot[parcol,thick,densely dotted,mark=square*,mark size=1.2pt] table[x=ki,y=chi2_f7_pure,col sep=comma]{plot_strategy3.csv};
\addlegendentry{$f=7$, pure CQD}
\addplot[parcol,thick,mark=square*,mark size=1.2pt] table[x=ki,y=chi2_f7_rlc,col sep=comma]{plot_strategy3.csv};
\addlegendentry{$f=7$, with $\RLC$}
\end{axis}
\end{tikzpicture}
\caption{Strategy 3: antiparallel set alone. Pure CQD never fits ($\chi^2\ge58$). With the retained-label
coefficient re-optimized at each $k_i$, $\chi^2$ stays near 1 over three decades for $f=7$, and for
$f=1$ it only bounds $k_i\lesssim1.5\times10^{-6}$.}
\label{fig:strategies}
\end{figure}

\paragraph{Why strategy 2 is preferred.} The parallel set is fitted by pure CQD with one parameter, so
$k_i$ is not influenced by any phenomenological term. Strategies 1 and 2 give the same $k_i$ within
0.5\%, which shows that the antiparallel data carry no independent information on $k_i$. Separating
the steps turns the antiparallel set into a clean test: with $k_i$ fixed, pure CQD has no freedom left
and fails ($\chi^2=109$), and the size of the retained-label fraction becomes a quantity to be
explained by a model of the field-reversal region.

%==========================================================================
\section{Dependence on the SG0 scaling \texorpdfstring{$f$}{f}}
\label{sec:fscan}

\begin{figure}[h]
\centering
\begin{tikzpicture}
\begin{groupplot}[group style={group size=2 by 1,horizontal sep=2cm},width=0.45\textwidth,height=0.37\textwidth,
  xlabel={SG0 field divisor $f$},grid=major,grid style={gray!20},xmin=0.5,xmax=10.5]
\nextgroupplot[ylabel={$k_i$},ymode=log,ymin=1e-7,ymax=1e-5]
\addplot[parcol,thick,mark=*,mark size=1.4pt] table[x=f,y=ki_par,col sep=comma]{plot_fscan.csv};
\nextgroupplot[ylabel={$\chi^2_\pp$ (15 points)},ymode=log,ymin=0.2,ymax=60]
\addplot[black,thick,mark=*,mark size=1.4pt] table[x=f,y=chi2_par,col sep=comma]{plot_fscan.csv};
\addplot[gray,dashed] coordinates{(0.5,23.7)(10.5,23.7)};
\end{groupplot}
\end{tikzpicture}
\caption{Parallel-only pure CQD fits as a function of the SG0 field divisor. Left: $k_i$. Right:
$\chi^2$; the dashed line is the 95\% limit for 14 degrees of freedom.}
\label{fig:fscan}
\end{figure}

The parallel fit is equally good for every $f\lesssim6$ ($\chi^2=0.46$--0.61), while $k_i$ grows from
$1.5\times10^{-7}$ to $1.8\times10^{-6}$. Beyond $f\approx6.5$ the fit degrades quickly:
$\chi^2=2.8$ at $f=7$, 11.9 at $f=8$ and 40 at $f=10$; the 95\% limit is crossed near $f\approx8.9$.
The data alone therefore measure $\alpha_0$, not $f$ and $k_i$ separately, and they cap $k_i$ at about
$3\times10^{-6}$ for any SG0 calibration. Between $f=6$ and $f=8$, $k_i$ changes from $1.8$ to
$2.9\times10^{-6}$, so the uncertainty of the factor 7 itself dominates the error on $k_i$.

\paragraph{Origin of the ceiling.} The initial slope of the parallel curve is
\begin{equation}
\frac{d\Delta D^\pp}{dI_0}\bigg|_0=\frac{13}{fI_\star}\;
\underbrace{\D_{c,1}\,\frac{\partial\Z^\pp/\partial\D_{c,0}}{\Z_0}}_{\equiv\,g(\D_{c,1})}.
\end{equation}
The capacity $g$ grows with $\D_{c,1}$ only until saturation sets in and peaks at $g\approx0.91$ near
$\D_{c,1}\approx2$ (Table~\ref{tab:g}). A measured slope therefore implies an upper bound on $f$.

\begin{table}[h]
\centering
\caption{Slope capacity $g(\D_{c,1})$ with $\DIR$ from the baseline.}
\label{tab:g}
\small\setlength{\tabcolsep}{4pt}
\begin{tabular}{lcccccccccc}
\toprule
$\D_{c,1}$ & 0.05 & 0.1 & 0.2 & 0.4 & 0.6 & 0.8 & 1.0 & 1.5 & 2.0 & 3.0\\
$\Z_0$ & 0.349 & 0.365 & 0.396 & 0.454 & 0.506 & 0.554 & 0.596 & 0.680 & 0.740 & 0.815\\
$g$ & 0.094 & 0.179 & 0.322 & 0.532 & 0.672 & 0.765 & 0.826 & 0.896 & 0.910 & 0.897\\
\bottomrule
\end{tabular}
\end{table}

\paragraph{Absolute geometry at $f=7$.} Solving $\DIR$ from the baseline gives $L_d=0.198$~m, which
contradicts both reported distances (0.320 and 0.403~m). Fixing $L_d=0.403$~m instead, the model
predicts a zero-current separation of 24.4~px rather than 13~px, which would require
$G_1\approx5.7$~T\,m$^{-1}$ instead of 10.8~T\,m$^{-1}$. The $f=7$ scaling brings $k_i$ close to the
numerical single-stage value ($2.51\times10^{-6}$) at the cost of moving the discrepancy into the
mechanical scale.

%==========================================================================
\section{Main caveats}
\begin{enumerate}
\item $f$, $\ell_0^{\rm eff}$ and $\ell_1^{\rm eff}$ are calibration inputs, not measured field integrals;
$k_i$ depends strongly on $f$ near $f=7$.
\item The absolute zero-current separation is taken as 13~px; it sets the vertical scale of every model
curve, since $\Delta D=13\,[\Z/\Z_0-1]$.
\item A single speed $v_a=600.5$~m\,s$^{-1}$ is used instead of the effusive speed distribution, and the
model computes branch centroids whereas the data are peak positions.
\item Reduced $\chi^2$ values of 0.04--0.2 indicate conservative or correlated uncertainties; a
generalized least-squares fit needs the covariance of the points.
\item The antiparallel fit depends on a phenomenological retained-label term that is not part of pure
CQD.
\end{enumerate}

%==========================================================================
\appendix
\section{Files}
All files are in the folder \texttt{fit\_20260706}.
\begin{description}
\item[\texttt{spin\_resolved\_sep\_20260706.csv}] \mbox{}\\ The data exported from the JLD2 file.
Columns: \texttt{I0\_A}, \texttt{dD\_parallel\_px},\\ \texttt{sigma\_parallel\_px},
\texttt{dD\_antiparallel\_px}, \texttt{sigma\_antiparallel\_px}.
\item[\texttt{fit\_curves\_f7\_20260706.csv}] \mbox{}\\ Recommended fit, $I_0=0$ to 3.70~A in steps of 0.01~A.
Columns: \texttt{I0\_A}, \texttt{dD\_parallel\_px}, \texttt{dD\_antiparallel\_px} (linear $w_{\rm ret}$),\\
\texttt{dD\_antiparallel\_exp\_px}, \texttt{dD\_antiparallel\_pureCQD\_px}, \texttt{Dc0}, \texttt{w\_ret}.
All values are changes relative to zero current (no 13-px offset).
\item[\texttt{fit\_curves\_f1\_20260706.csv}] \mbox{}\\ The same for $f=1$.
\item[\texttt{fit\_params\_f7\_20260706.csv}] \mbox{}\\ Parameters of the recommended fit.
\item[\texttt{results.json}] All fit results, including strategies 1--4 and the $f$ scan.
\item[\texttt{cqd\_model.js}, \texttt{run\_fits.js}] \mbox{}\\ Model and fits; \texttt{node run\_fits.js}
regenerates every number and CSV file (about 20~s). \texttt{make\_table.js} and
\texttt{build\_report.js} regenerate Table~\ref{tab:points} and this document.
\end{description}

\end{document}
